% !TEX program = xelatex
\documentclass{英语试卷}
\begin{document}
\examtitle{南京市、盐城市 2025 届高三年级第一次模拟考试}{英语}
\par 南京市、盐城市 2025 届高三年级第一次模拟考试\par
\par 英语试题卷\par
\begin{notice}
\par 1. 答卷前，考生务必将自己的姓名、准考证号填写在答题卡上。\par
\par 2. 回答选择题时，选出每小题答案后，用铅笔把答题卡上对应题目的答案标号涂黑。如需改动，用橡皮擦干净后，再选涂其他答案标号。回答非选择题时，将答案写在答题卡上，写在本试卷上无效。\par
\par 3. 考试结束后，将本试卷和答题卡一并交回，\par
\end{notice}
\parttitle{第一部分 听力(共两节，满分 30 分)}
\par 做题时，先将答案标在试卷上。录音内容结束后，你将有两分钟的时间将试卷上的答案转涂到答题卡上。\par
\sectitle{第一节(共 5 小题；每小题 1.5 分，满分 7.5 分)}
\par 听下面 5 段对话，每段对话后有一个小题，从题中所给的 A、B、C 三个选项中选出最佳选项。听完每段对话后，你都有 10 秒钟的时间来回答有关小题和阅读下一小题，每段对话仅读一遍。\par
\question{1. What are the speakers mainly talking about?}
\begin{choices}(3)
  \item The woman's illness.
  \item A maths lecture.
  \item The school library.
\end{choices}
\question{2. When will John see the doctor on the 13th ?}
\begin{choices}(3)
  \item At 10:45 a. m.
  \item At 1:30 p. m.
  \item At 2:00 p. m.
\end{choices}
\question{3. What does the woman mean?}
\begin{choices}(3)
  \item She is treated unfairly.
  \item She is tired of working.
  \item She is highly valued.
\end{choices}
\question{4. What mistake did the man make?}
\begin{choices}(3)
  \item He misspelled the address.
  \item He downloaded a wrong file.
  \item He left out a letter.
\end{choices}
\question{5. What does the man think of the hat?}
\begin{choices}(3)
  \item It's a good buy.
  \item It's not worth it.
  \item It's out of fashion.
\end{choices}
\sectitle{第二节(共 15 小题;每小题 1.5 分,满分 22.5 分)}
\par 听下面 5 段对话或读白，每段对话或独白后有几个小题，从题中所给的 A、B、C 三个选项中选出最佳选项。听每段对话或独白前，你将有时间阅读各个小题，每小题 5 秒钟。听完后，各小题将给出 5 秒钟的作答时间。每段对话或独白读两遍。\par
\par 听第 6 段材料, 回答题 6、7 题。\par
\question{6. What is the woman's job?}
\begin{choices}(3)
  \item A meeting organizer.
  \item An interpreter.
  \item A secretary.
\end{choices}
\question{7. What does the woman request the man to do?}
\begin{choices}(3)
  \item Use shortened sentences.
  \item Make the discussion lively.
  \item Allow certain breaks.
\end{choices}
\question{8. What happened to the missing man when he reappeared?}
\begin{choices}(3)
  \item He had an Australian passport.
  \item He had a different appearance.
  \item He had a new name.
\end{choices}
\question{9. What revealed the missing man's true identity?}
\begin{choices}(3)
  \item The clothes on the beach.
  \item A label on his suitcase.
  \item His check-in record.
\end{choices}
\question{10. How does the woman feel after learning the story?}
\begin{choices}(3)
  \item Interested.
  \item Confused.
  \item Bored.
\end{choices}
\par 听第 8 段材料，回答第 11 至 13 题。\par
\question{11. What is Jack worried about?}
\begin{choices}(3)
  \item His tight holiday schedule.
  \item A fine for his overdue book.
  \item No access to the library website.
\end{choices}
\question{12. What could Jack do to extend the period?}
\begin{choices}(3)
  \item Click the "Borrow" button.
  \item Visit the school library in person.
  \item Operate online under his account.
\end{choices}
\question{13. What does Emily plan to do?}
\begin{choices}(3)
  \item Watch the Dune film.
  \item Help Jack return the book.
  \item Read Dune after Jack finishes it.
\end{choices}
\par 听第 9 段材料，回答第 14 至 17 题。\par
\question{14. What is the probable relationship between the speakers?}
\begin{choices}(3)
  \item Coach and trainee.
  \item Husband and wife.
  \item Fellow workers.
\end{choices}
\question{15. Why is the advanced class so hard for the woman?}
\begin{choices}(3)
  \item She is lacking in flexibility.
  \item The instructor is hard on her.
  \item There is no rest between exercises.
\end{choices}
\question{16. What does the man think of the advanced class?}
\begin{choices}(3)
  \item Unchallenging.
  \item Amusing.
  \item Tiring.
\end{choices}
\question{17. What is the man probably going to do first tomorrow?}
\begin{choices}(3)
  \item Lift weights.
  \item Prepare a cake.
  \item Join the woman.
\end{choices}
\par 听第 10 段材料, 回答第 18 至 20 题。\par
\question{18. What makes small talk challenging?}
\begin{choices}(3)
  \item No time for preparation.
  \item Too many facts to check.
  \item Heavy reliance on bonding.
\end{choices}
\question{19. What are small talk participants expected to do according to Rachel?}
\begin{choices}(3)
  \item To give quick responses.
  \item To show interest in the talk.
  \item To speak in a fascinating way.
\end{choices}
\question{20. What will the speaker talk about next?}
\begin{choices}(3)
  \item Importance of small talk skills.
  \item Consequences of quick reaction.
  \item Tips on pacing yourself in small talk.
\end{choices}
\parttitle{第二部分 阅读理解（共两节，满分 50 分）}
\sectitle{第一节（共15小题；每小题2.5分，满分37.5分）}
\par 阅读下列短文，从每题所给的 A、B、C 和 D 四个选项中，选出最佳答案，\par
\passagelabel{A}
\par This coffee cup is a high-quality product designed for coffee lovers. It is made with great attention to detail to enhance your coffee-drinking experience.\par
\par The cup is made of high-quality ceramic. Ceramic is known for its ability to maintain heat well, keeping your coffee warm for a longer period. It also gives a smooth and pleasant touch when you hold the cup. The handle is firmly attached to the cup. It is made of the same ceramic material, ensuring durability and a seamless (无缝的) look.\par
\par [图片]\par
\par Design\par
\par ◆ Th coffee cup has a standard capacity of 300ml. This size is perfect for a regular cup of coffee, whether it is\par
\par a latte (拿铁), espresso, or cappuccino.\par
\par ◆ Th cup has a classic round shape. The opening is wide enough for easy sipping and also allows the pleasant smell to be fully released.\par
\par Usage Instructions\par
\par Before using the coffee cup for the first time, it is recommended to wash it with warm water and mild cleaner.\par
\par When pouring coffee into the cup, be careful not to overfill it.\par
\par ◆ After use, wash the cup instantly. It is dishwasher-safe, but hand-washing can also be used to extend its lifespan (使用寿命).\par
\par Maintenance\par
\par ◆ Avoid sudden temperature changes. For example, do not pour extremely hot coffee into a cold cup or put a hot cup directly into cold water, as this may cause the ceramic to break.\par
\par ◆ Store the coffee cup in a dry and clean place.\par
\question{21. Which section can readers refer to for storage requirements?}
\begin{choices}(2)
  \item Material.
  \item Design.
  \item Usage Instructions.
  \item Maintenance.
\end{choices}
\question{22. What do we know about the cup's material?}
\begin{choices}(2)
  \item It is sensitive to mild cleaner.
  \item It has properties of absorbing heat.
  \item It helps smooth the holder's hand skin.
  \item It is subject to sudden temperature changes.
\end{choices}
\question{23. What is the text?}
\begin{choices}(2)
  \item A product report.
  \item A user guide.
  \item A cup advertisement.
  \item A coffee recipe
\end{choices}
\passagelabel{B}
\par Inspired by an event held by the Dallas-based Daddy Stroller (婴儿车) Social Club (DSSC), a group for fathers, Brandon Shaw wrote an 83-page typed letter to his father. From the Blacklit bookstore, Kalvin Bridgewater led 15 dads, including his own father, who was seated next to Shaw, in an emotional event where the men dug into what they hated, loved, misunderstood and appreciated about their fathers as they grew up.\par
\par Founded by Bridgewater in 2021, DSSC initially met with about 30 dads per event but skyrocketed in popularity after two online videos became hits in April 2024. It now attracts over 150 dads at events, where fathers bond over shared experiences and provide emotional backing for one another.\par
\par DSSC's appeal is broad, attracting dads from various professions, including doctors, engineers, and electricians. The group emphasizes emotional support, with members regularly checking in on one another. As Shaw shared, the group provides a space where men can be open, asking, "How are you feeling?" Shaw, who joined the group when he first became a father, said it has made him more patient and understanding.\par
\par Bridgewater, who became a certified doula (产妇陪护师) to better support his wife through pregnancy (孕期), also ensures DSSC allows partners some time for themselves. At a May event, scheduled just before Mother's Day, mothers had the day off while fathers spent time bonding with their children.\par
\par Brandon Shaw described the experience as transformative, saying it helped him connect with his own weaknesses, something he had previously kept hidden. He sent the letter he wrote to his father for Father's Day. It arrives soon. He said he doesn't expect a response and just hopes his father reads it. As Bridgewater puts it, fatherhood is just about "being honest and sharing".\par
\question{24. What can be inferred about DSSC from the first two paragraphs?}
\begin{choices}(2)
  \item It creates a group of father influencers.
  \item It repairs broken father-son relationships.
  \item It promotes fathers' emotional well-being.
  \item It highlights fathers' vital role in childcare.
\end{choices}
\question{25. What makes DSSC appealing?}
\begin{choices}(2)
  \item Parenting instructions.
  \item Good childcare service.
  \item Supportive atmosphere
  \item Diverse member backgrounds.
\end{choices}
\question{26. Why is the May event mentioned in paragraph 4?}
\begin{choices}(2)
  \item To illustrate another aim of DSSC.
  \item To justify the founder's personal experience.
  \item To stress the importance of DSSC.
  \item To show an ideal husband-wife relationship.
\end{choices}
\question{27. What is the main idea of the last paragraph?}
\begin{choices}(2)
  \item The severity of parental issues.
  \item DSSC's influence on its members.
  \item The necessity to identify weaknesses.
  \item Bridgewater's devotion to fatherhood.
\end{choices}
\passagelabel{C}
\par Dancing lowers the depression associated with Parkinson's disease, and the benefits can be seen in multiple ways, a new study shows.\par
\par "It was very cool to see dance had a positive effect on the mood circuits in the brain, which we could see in the imaging from MRI brain scans (扫描)," says Joseph DeSouza, author of the study. "These improvements were also reported by the participants via survey. Our study is the first to demonstrate these benefits across these two detection methods."\par
\par The study followed 23 participants with Parkinson's disease from the Sharing Dance Parkinson's program at Canada's National Ballet School, along with family members or caregivers. Participants took weekly dance classes for 8 months, progressing from simple leg and foot work to more complex movements like waltzes and more. The research team focused on the subcallosal cingulate gyrus (SCG), a brain region linked to depression. They measured mood and depression before and after each class using the Geriatric Depression Scale (GDS) (老年抑郁量表) and conducted regular MRI scans.\par
\par The researchers found that after each dance class, reported depression rates dropped, and the effect was increasing from class to class, with significant improvements after 8 months. "We showed that SCG's BOLD (blood oxygen level dependent) signal decreases while dancing, which means the SCG was not functioning as fast as it would if you had depression," says DeSouza. "Parkinson’s disease is often accompanied by depression long before motor symptoms (症状) appear."\par
\par Dance is thought to have a double benefit, with music activating the brain's reward centers, and the movement acting on sensory and motor circuits. DeSouza, dancing with participants for 14 years, says that while dancing is not a treatment for Parkinson's disease by itself, the benefits are clear. "We're not trying to cure it with dance," he says. "What we're trying to do is to have people live a better quality of life. This goes for both the patients, and their families that take care of them."\par
\question{28. What do we know about the results of the two detection methods?}
\begin{choices}(2)
  \item They fall short of expectations.
  \item They are more accurate than before.
  \item They are distinct from each other.
  \item They show consistent conclusions.
\end{choices}
\question{29. How did the researchers conduct the study?}
\begin{choices}(2)
  \item By referring to previous studies.
  \item By making pre-post comparisons.
  \item By employing computer modeling.
  \item By conducting drug experiments.
\end{choices}
\question{30. What did the researchers find about the SCG?}
\begin{choices}(2)
  \item It is in charge of motor system.
  \item It can be measured by the GDS.
  \item Dancing slows down its functioning.
  \item Depression weakens its BOLD signal.
\end{choices}
\question{31. What is the best title of the text?}
\begin{choices}(2)
  \item Dancing Away Parkinson's Blues
  \item Dancing: Multiple Benefits To Be Seen
  \item Dancing With Your Mood Circuits
  \item Dancing: A Journey To Mental Well-being
\end{choices}
\passagelabel{D}
\par Imagine an artificial intelligence model that can watch and understand moving images like you and me. Now, scientists have made this a reality by creating MovieNet: an innovative AI that processes videos much like how our brains interpret real-life scenes.\par
\par This brain-inspired AI model can perceive moving scenes by simulating how neurons (神经元) make real-time sense of the world. While conventional AI excels at still image recognition, MovieNet is designed to understand complex, dynamic scenes, allowing it to recognize small changes over time.\par
\par To create MovieNet, Hollis Cline, senior author of the study, examined how animals' brains process real-world scenes. Specifically, he studied how tadpoles' (蝌蚪的) neurons respond to moving stimuli. Tadpoles' neurons process these visual inputs like solving puzzle pieces: each piece alone may not make sense, but together they form a complete image in motion. Different neurons process various "puzzle pieces" of a real-life moving image, which the brain then combines to form a continuous scene.\par
\par Using this understanding, the researchers trained MovieNet to process video clips by identifying small visual signals, allowing it to distinguish between small changes in scenes. In testing, MovieNet achieved 82.3\% accuracy, outperforming trained human observers by 18\% and doing much better than existing AI models.\par
\par Beyond its higher accuracy, MovieNet is more environment-friendly. Conventional AI requires vast amounts of data and energy, but MovieNet compresses (压缩) data like a zipped file, reducing environmental impact. "We’ve also made AI far less demanding," says Cline. MovieNet could also transform medicine. It might help detect early signs of diseases by tracking small movement changes that are difficult for human eyes to detect.\par
\par Looking ahead, the team plans to improve MovieNet's ability to adapt to different environments and enhance its potential applications. "By designing models that think like living organisms (有机体), we can achieve efficiency levels that aren’t possible with traditional AI," says Cline, "Biology is a gold mine."\par
\question{32. What does the underlined word "simulating" mean in paragraph 2?}
\begin{choices}(2)
  \item Copying.
  \item Inspiring.
  \item Demonstrating.
  \item Recognizing.
\end{choices}
\question{33. MovieNet is trained \blanknum{}.}
\begin{choices}(2)
  \item to integrate dynamic images
  \item to increase visual inputs
  \item to feed on information pieces
  \item to look into tadpoles' neurons
\end{choices}
\question{34. According to paragraph 5, what is an advantage of MovieNet?}
\begin{choices}(2)
  \item Its powerful data storage.
  \item Its low energy consumption.
  \item Its control of medical diseases.
  \item Its improvement on drug effects.
\end{choices}
\question{35. What can be concluded from Cline's words in the last paragraph?}
\begin{choices}(2)
  \item Biology brings us financial profits.
  \item Living organisms can outperform AI.
  \item AI can assist in biological development.
  \item Drawing inspiration from biology is promising.
\end{choices}
\sectitle{第二节（共5小题；每小题2分，满分12.5分）}
\par 根据短文内容，从短文后的选项中选出能填入空白处的最佳选项。选项中有两项为多余选项。\par
\par Growing up, I found the word Saturday almost equaled cleaning in my family. As soon as any child was able to use a dust cloth, he or she was assigned a chore (家务活). \blanknum{36} Exceptions were made of course, but the older we grew, the more my parents expected us to balance these responsibilities.\par
\par Chores were part of my family's way of functioning. We complained and tried to avoid them, but we never questioned why we did them. As a family, we all had roles in keeping the home running smoothly, and the responsibility grew with age. \blanknum{37} 微信公众号 IAI English\par
\par When I was a kid, I knew when doing chores I was helping my family. It wasn't until I grew into adulthood that I realized that chores are more than a list of tasks. \blanknum{38} Being assigned chores communicates that each family member is needed and that their action or non-action affects every single other person living under the same roof.\par
\par Handling daily household chores is never easy, but doing them out of love for my family makes it manageable. \blanknum{39} I enjoy jumping in to help whenever my parents need a hand. While working together, we have the chance to have meaningful conversations, share experiences, and build memories.\par
\par \blanknum{40} Household chores mean learning how to manage responsibility and understand kids' role, not to mention practicing discipline, time management, and selflessness. All these consequences would seem to tip the scales (天平) in favor of chores.\par
\begin{choices}(1)
  \item It made sense to us, as much as we disliked it.
  \item They create a sense of belonging within a family.
  \item Household chores were a source of great stress for us.
  \item Actually, doing chores alongside my family can make them a joy.
  \item We had to complete our assigned chore before school assignments.
  \item Allowing children to share burdens helps them recognize their strengths.
  \item It also comes as no surprise that chores are important for personal growth.
\end{choices}
\parttitle{第三部分 语言知识运用（共两节，满分 30 分）}
\sectitle{第一节（共15小题；每小题1分，满分15分）}
\par 阅读下面短文，从短文后各题所给的A、B、C和D四个选项中，选出可以填入空白处的最佳选项，\par
\par On a fall day, my wife and I went boating on the Blackfoot River. I was excited about the trip because it would allow me to forget my \blanknum{41}.\par
\par At noon, we rested on a bank. \blanknum{42} the scenery, we barely realized the approaching danger. Then three large \blanknum{43} emerged from the forest: a huge mama bear with two babies. They were less than 100 feet away across the river, far closer than the recommended \blanknum{44} distance.\par
\par I could feel \blanknum{45} reached deep into my stomach. Time stood still. Then, I grabbed the bear spray (喷雾) and we walked downriver to \blanknum{46} ourselves. "Hey, bear!" I said, "We are not here to \blanknum{47} you." At that moment, I thought the mama bear would \blanknum{48} across the river, covering the 100 feet in seconds. Although I knew how to use the spray, I didn't think I could use it \blanknum{49}.\par
\par To our surprise, the mama bear \blanknum{50} drank water from the river. When one baby wandered just a few steps away, she quickly \blanknum{51} him. Just like their \blanknum{52} appearance, they were gone, disappearing into the forest. Ne returned to our boat, lost in \blanknum{53}.\par
\par This encounter \blanknum{54} me, making my daily concerns seem insignificant in the larger context of our planet. The experience, the fear, and my \blanknum{55} as a human have changed my perspective on the world forever.\par
\cz{41}{concerns}{principles}{passions}{secrets}
\cz{42}{Tired of}{Absorbed in}{Looking into}{Reflecting on}
\cz{43}{signs}{pictures}{targets}{shapes}
\cz{44}{social}{safe}{average}{constant}
\cz{45}{fear}{surprise}{thrill}{courage}
\cz{46}{relax}{corner}{dismiss}{distance}
\cz{47}{bathe}{tease}{bother}{pet}
\cz{48}{sail}{float}{charge}{wander}
\cz{49}{in advance}{at random}{at peace}{in time}
\cz{50}{probably}{simply}{naturally}{certainly}
\cz{51}{rescued}{abandoned}{fetched}{examined}
\cz{52}{timely}{inappropriate}{unexpected}{regular}
\cz{53}{awe}{shock}{regret}{calmness}
\cz{54}{astonished}{humbled}{refreshed}{fueled}
\begin{choices}(3)
  \item limitations 3. responsibilities
  \item advantages
  \item virtues
\end{choices}
\sectitle{第二节（共10小题；每小题1.5分，满分15分）}
\par 阅读下面短文，在空白处填入 1 个适当的单词或括号内单词的正确形式，\par
\par In the digital age, many young people tend to share marriage announcements on social media platforms like Sina Weibo. Yet a new trend has reemerged — publishing marriage announcements in traditional newspapers. \blanknum{56} quick search on the topic "publishing marriage announcements in newspapers" will bring up \blanknum{57}(variety) posts, along with those who have done so sharing their experiences.\par
\par \blanknum{58} short-lived digital posts, newspaper announcements provide a sense of permanence, offering a meaningful way \blanknum{59} (keep) memories during a time when digital content can't last. Some young people view publishing marriage announcements in newspapers as more formal than \blanknum{60} in social media, and they describe it as a \blanknum{61} (commit) written in black ink and on white paper — a declaration of the couple's promise to grow old together.\par
\par However, this practice is not a new phenomenon in China. In the first half of the twentieth century, when marriage laws were still underdeveloped, such announcements \blanknum{62} (carry) legal significance. They allowed couples to \blanknum{63} (public) share their joyous news with their family and friends.\par
\par The regained popularity serves as an example of \blanknum{64} tradition can better adapt to modern times. By combining cultural heritage with personal creativity, it offers meaningful rituals (仪式) in today's digital world. These announcements, as symbols of love and individuality, remind us that some traditions are worth \blanknum{65} (preserve).\par
\parttitle{第四部分 写作（共两节，满分 40 分）}
\sectitle{第一节（满分15分）}
\par 假定你是李华，外教 Mr. Chuck 发现你的作文完全由 AI 生成，判为不及格。请你给外教写一封邮件，内容包括：\par
\par (1) 承认错误;\par
\par (2) 做出许诺;\par
\par (3) 补救措施。 注意：\par
\par (1) 写作词数应为 80 个左右;\par
\par (2) 请按如下格式在答题卡的相应位置作答:\par
\par Dear Mr. Chuck,\par
\sectitle{第二节（满分25分）}
\par 阅读下面材料，根据其内容和所给段落开头语续写两段，使之构成一篇完整的短文。\par
\par After my family moved to America, we lived on Pearl Street in New York's Chinatown. Below our flat, my family ran a bakery. My mother named me after the street: Pearl Wong.\par
\par A typical Chinese-style mum, my mother often reminded me of the importance of being outstanding. "Be number one, or be nothing," she demanded in broken English. She would check my report cards and frown at anything lower than an A.\par
\par At a holiday party, my brother received a chess set, missing a pawn (卒) and a knight(马). My mother dismissed it as unwanted junk, but we were fascinated, especially me.\par
\par "Why can't they move more steps?" "Why do they go crossways to take other men?" I had so many questions. My mother patted the flour off her hands. "Let me see book," she said. "Ah. Don't know why? You find out yourself."\par
\par So I read and memorized the rules, studied the moves, and practiced endlessly. Soon, I was beating my brothers. With every game, my chess improved, my wins grew, and my hunger for victory deepened.\par
\par I started entering tournaments and winning trophies (奖杯). My mother became my biggest supporter — but in her own way. She pushed me to win every game and displayed every trophy in the bakery window. I asked her several times not to do so but she insisted. Soon, the cakes were surrounded by rows of shining trophies. "This my daughter, Pearl," she would tell others, her voice full of pride. But the more she showed me off, the more I felt trapped, as if I were no longer playing for myself, as if I were her trophy.\par
\par One Saturday in Chinatown, my mother paraded (招摇过市) me through the stores, without making a purchase. "This my daughter, Pearl Wong, Chinatown chess champion," she announced. The shopkeepers, with forced smiles, looked at us with an obviously growing impatience. 注意：\par
\par （1）续写词数应为 150 左右；\par
\par （2）请按如下格式在答题卡的相应位置作答。\par
\writingpara{Paragraph 1: Yet, she went on and on, totally ignoring my burning eyes.}
\writinglines{8}
\writingpara{Paragraph 2: When I finally returned to the bakery, my mother was putting away the trophies. \blanknum{}}
\writinglines{8}
\end{document}
